\thispagestyle{fancy}
\fancyhead[R]{A.U 2025-2026}
\fancyhead[L]{Chapitre 2 : Analyse et spécification des besoins}

\chapter{Analyse et spécification des besoins}

\section*{Introduction}
Dans le premier chapitre, nous avons choisi d'adopter la méthodologie Scrum pour
concevoir notre futur système. Dans ce chapitre, nous procédons à une analyse détaillée
des besoins liés au développement de notre plateforme. Nous commençons par identifier
les acteurs ainsi que les besoins fonctionnels et non fonctionnels. Par la suite, nous
présentons le backlog produit et la planification des sprints. Enfin, nous exposons
l'environnement technique retenu pour le développement, ainsi que l'architecture
générale de l'application.

\section{Modélisation du contexte}
Dans le cycle de vie du développement logiciel, la spécification des besoins fonctionnels
et non fonctionnels représente la première étape de ce cycle. Elle permet de spécifier
les attentes de la plateforme à concevoir, tout en tenant compte d'une compréhension
claire des besoins de l'organisation.

\subsection{Identification des acteurs}
La plateforme de supervision met en interaction plusieurs acteurs qui participent à la
surveillance et à l'administration de l'infrastructure serveur. Chaque acteur dispose
d'un rôle spécifique afin d'assurer le bon fonctionnement du système.

\begin{itemize}[label=\ding{109}]
\item \textbf{Administrateur :} utilisateur authentifié de la plateforme, via
l'application web ou mobile. Il supervise l'état des serveurs, consulte les métriques en
temps réel, gère les alertes, exécute les actions de maintenance à distance (redémarrage
ou arrêt de services et de serveurs), et supervise la chaîne d'intégration continue et
l'état du cluster Kubernetes.

\item \textbf{Agent de supervision :} acteur automatisé, matérialisé par un programme
Python exécuté de manière autonome sur chaque serveur supervisé. Il collecte
périodiquement les métriques système (CPU, RAM, disque, réseau, uptime) ainsi que la
liste des services actifs toutes les 5 secondes, puis les transmet au backend via l'API
REST, sans intervention humaine.
\end{itemize}

\subsection{Spécification des besoins}
Notre projet a pour but l'analyse, la conception et la réalisation d'une plateforme de
supervision et de gestion à distance d'infrastructure serveur. La solution développée
doit répondre aux exigences et aux besoins déterminés tout au long de la phase d'étude.

\subsubsection*{Besoins fonctionnels}
Les exigences fonctionnelles représentent l'ensemble des tâches et des fonctionnalités
que le système doit assurer afin de répondre aux besoins des utilisateurs. Elles sont
présentées, dans le tableau~\ref{tab:besoins-fonctionnels}, en fonction de chaque acteur
de la plateforme.

\renewcommand{\arraystretch}{1.5}
\begin{longtable}{|>{\raggedright\arraybackslash}p{3.8cm}
                  |>{\raggedright\arraybackslash}p{10.4cm}|}
\hline
\textbf{Acteurs} & \textbf{Fonctionnalités} \\
\hline
\endfirsthead

\hline
\textbf{Acteurs} & \textbf{Fonctionnalités} \\
\hline
\endhead

\hline
\multicolumn{2}{r}{\textit{Suite du tableau à la page suivante}} \\
\endfoot

\hline
\caption{Besoins fonctionnels par acteur}
\label{tab:besoins-fonctionnels}
\endlastfoot

Administrateur &
\textit{Supervision et administration :}
\begin{itemize}[label=--, leftmargin=*, nosep, topsep=2pt]
  \item S'authentifier à la plateforme (web et mobile)
  \item Consulter le tableau de bord global des serveurs
  \item Visualiser les métriques en temps réel (CPU, RAM, disque, réseau)
  \item Consulter l'historique des métriques sous forme de graphiques
  \item Consulter la liste des alertes actives
  \item Acquitter une alerte prise en charge
  \item Recevoir des notifications par courrier électronique
  \item Redémarrer ou arrêter un service applicatif à distance
  \item Redémarrer complètement un serveur à distance
  \item Consulter la liste des services détectés sur un serveur
  \item Consulter l'historique des sauvegardes
  \item Consulter le journal d'audit des actions effectuées
\end{itemize}
\vspace{4pt}
\textit{Industrialisation et déploiement :}
\begin{itemize}[label=--, leftmargin=*, nosep, topsep=2pt]
  \item Déclencher et suivre l'exécution du pipeline CI/CD
  \item Consulter les rapports d'analyse de qualité du code
  \item Consulter les rapports de scan de vulnérabilités
  \item Superviser l'état du cluster Kubernetes (pods, ressources)
  \item Gérer les secrets de l'infrastructure
  \item Consulter les tableaux de bord de supervision d'infrastructure
\end{itemize} \\
\hline

Agent de supervision &
\begin{itemize}[label=--, leftmargin=*, nosep, topsep=2pt]
  \item Collecter les métriques système toutes les 5 secondes
  \item Détecter automatiquement les services actifs du serveur
  \item Transmettre les données collectées au backend via l'API REST
  \item Poursuivre la collecte en cas d'indisponibilité du backend
  \item Reprendre l'envoi dès le rétablissement de la connexion
\end{itemize} \\
\hline

\end{longtable}

\subsubsection*{Besoins non fonctionnels}
Les besoins non fonctionnels décrivent les contraintes techniques, les critères de
qualité ainsi que les exigences de performance auxquelles la plateforme doit répondre
afin de garantir son bon fonctionnement.

\begin{itemize}[label=\ding{220}]
\item \textbf{Sécurité :}
\begin{itemize}[label=\ding{43}]
    \item Authentifier les utilisateurs à l'aide d'un identifiant et d'un mot de passe
          stocké sous forme hachée ;
    \item Sécuriser les échanges entre le frontend et le backend grâce à des jetons JWT
          d'une durée de validité limitée ;
    \item Protéger les actions sensibles par une double vérification : validité du jeton
          et appartenance au rôle administrateur, contrôlée en base de données ;
    \item Journaliser l'ensemble des actions d'administration à distance dans un journal
          d'audit.
\end{itemize}

\item \textbf{Performance :}
\begin{itemize}[label=\ding{43}]
    \item Assurer une communication en temps réel entre le serveur et les clients grâce
          au protocole WebSocket ;
    \item Garantir un délai de rafraîchissement de l'interface inférieur à une seconde ;
    \item Optimiser les requêtes de consultation des métriques par l'indexation de la
          base de données.
\end{itemize}

\item \textbf{Disponibilité et maintenabilité :}
\begin{itemize}[label=\ding{43}]
    \item Garantir le redémarrage automatique des conteneurs en cas de défaillance grâce
          à l'orchestrateur Kubernetes ;
    \item Structurer le code source du backend selon une architecture modulaire (routes,
          modèles, services, middlewares) ;
    \item Conteneuriser l'ensemble des composants avec Docker afin de faciliter le
          déploiement et la maintenance.
\end{itemize}

\item \textbf{Ergonomie :}
\begin{itemize}[label=\ding{43}]
    \item Proposer une interface web responsive et intuitive ;
    \item Assurer une cohérence visuelle entre l'application web et l'application mobile ;
    \item Présenter l'ensemble de l'interface en langue française.
\end{itemize}

\item \textbf{Interopérabilité :}
\begin{itemize}[label=\ding{43}]
    \item Permettre la communication entre les clients web et mobile et le backend à
          travers une API REST unique ;
    \item Assurer l'échange de données entre l'agent Python et le backend via le
          protocole HTTP ;
    \item Établir des connexions SSH vers les serveurs supervisés pour l'exécution des
          actions à distance.
\end{itemize}

\item \textbf{Scalabilité :}
\begin{itemize}[label=\ding{43}]
    \item Supporter l'ajout de nouveaux serveurs supervisés sans modification du backend,
          par simple installation de l'agent ;
    \item Séparer les couches présentation, métier et données afin de permettre leur
          évolution indépendante ;
    \item Assurer une gestion flexible et évolutive des données grâce à MongoDB, adapté
          au stockage de volumes importants de métriques.
\end{itemize}
\end{itemize}

\subsection{Diagramme de cas d'utilisation général}
Dans cette section, nous décrivons les exigences du système de manière formelle en
utilisant le diagramme de cas d'utilisation global. La figure~\ref{fig:usecase-global}
présente l'ensemble des fonctions fournies par l'application. Tous les cas d'utilisation
nécessitent une authentification comme condition préalable.

\begin{figure}[H]
    \centering
    \includegraphics[width=1\textwidth]{Image/UseCaseGlobal.png}
    \caption{Diagramme de cas d'utilisation global}
    \label{fig:usecase-global}
\end{figure}

\section{Analyse globale}
Nous présentons, dans ce qui suit, une analyse globale comportant le diagramme de
classes ainsi que l'architecture de la solution.

\subsection{Diagramme de classes d'analyse}
Bien que l'application repose sur une base de données NoSQL~\cite{mongodb}, dont le
modèle de stockage est fondé sur des \textbf{collections} et des \textbf{documents}
plutôt que sur des tables relationnelles, il demeure pertinent d'élaborer un diagramme
de classes d'analyse afin de représenter les principales entités du système. Cette
modélisation conceptuelle ne vise pas à reproduire fidèlement la structure physique de
la base de données, mais à offrir une vision abstraite des différents objets métier
manipulés par l'application.

Elle permet ainsi d'identifier les principales collections, leurs attributs essentiels
ainsi que les relations fonctionnelles qui les unissent, facilitant la compréhension de
l'architecture des données et des interactions entre les différents composants du
système. La figure~\ref{fig:diagramme-classes} présente ce diagramme.

\begin{figure}[H]
    \centering
    \includegraphics[width=1\textwidth]{Image/DiagrammeClasses.png}
    \caption{Diagramme de classes d'analyse global}
    \label{fig:diagramme-classes}
\end{figure}

\subsection{Architecture globale}
La définition de l'architecture logicielle constitue une étape essentielle dans la
conception de notre plateforme. En effet, le choix de l'architecture influence
directement les performances, la maintenabilité, l'évolutivité ainsi que la capacité du
système à satisfaire les exigences fonctionnelles et non fonctionnelles identifiées lors
de l'analyse des besoins.

Dans cette perspective, une étude comparative a été menée entre trois architectures
largement adoptées dans le développement des applications modernes : l'architecture
monolithique, l'architecture client--serveur à trois niveaux (N-Tiers) et l'architecture
orientée microservices. Cette comparaison, présentée dans le
tableau~\ref{tab:comparaison-archi}, vise à mettre en évidence les caractéristiques, les
avantages et les limites de chacune de ces approches.

\begin{longtable}{|p{3.3cm}|>{\raggedright\arraybackslash}p{10.9cm}|}
\hline
\textbf{Architecture} & \textbf{Caractéristiques} \\
\hline
\endfirsthead

\hline
\textbf{Architecture} & \textbf{Caractéristiques} \\
\hline
\endhead

\hline
\multicolumn{2}{r}{\textit{Suite du tableau à la page suivante}} \\
\endfoot

\hline
\caption{Comparaison des architectures Monolithique, N-Tiers et Microservices}
\label{tab:comparaison-archi}
\endlastfoot

\textbf{Monolithique} &
\begin{itemize}[leftmargin=*, nosep]
  \item L'application est développée et déployée sous la forme d'un unique bloc logiciel.
  \item Les différents modules sont fortement couplés.
  \item Le déploiement est relativement simple.
  \item Toute évolution nécessite le redéploiement complet de l'application.
  \item La maintenance et l'évolutivité deviennent complexes lorsque l'application
        grandit.
\end{itemize} \\
\hline

\textbf{N-Tiers} &
\begin{itemize}[leftmargin=*, nosep]
  \item Séparation claire des responsabilités entre les couches de présentation, métier
        et données.
  \item Déploiement indépendant du frontend et du backend.
  \item Base de données centralisée partagée par les différentes couches.
  \item Facilite les opérations de maintenance, de test et d'évolution.
  \item Offre un bon compromis entre modularité, simplicité de développement et
        performances.
\end{itemize} \\
\hline

\textbf{Microservices} &
\begin{itemize}[leftmargin=*, nosep]
  \item Décomposition de l'application en services indépendants et faiblement couplés.
  \item Chaque service dispose généralement de sa propre base de données.
  \item Permet un déploiement et une mise à l'échelle indépendants de chaque service.
  \item Les échanges entre services s'effectuent via des API REST.
  \item Offre une grande flexibilité mais implique une orchestration plus complexe.
\end{itemize} \\
\hline

\end{longtable}

Au terme de cette étude comparative, l'architecture client--serveur à trois niveaux
(\textbf{N-Tiers}) a été retenue pour la conception de notre plateforme. Ce choix
s'explique par sa capacité à assurer une séparation claire des responsabilités entre les
différentes couches de l'application, tout en garantissant une bonne maintenabilité, une
évolutivité satisfaisante et une mise en œuvre relativement simple. Contrairement à
l'architecture monolithique, elle favorise une meilleure modularité, tandis qu'elle
évite la complexité de déploiement et d'orchestration inhérente à une architecture
basée sur les microservices. Elle constitue ainsi une solution adaptée aux besoins
techniques et fonctionnels du projet.

\subsubsection{Architecture logique de la plateforme}
L'architecture logicielle joue un rôle essentiel dans le développement d'un système
fiable, efficace et capable d'évoluer en fonction des besoins. Dans le cadre de ce
projet, nous avons opté pour une architecture logique organisée en plusieurs couches
afin de garantir une répartition claire des responsabilités entre les différents
composants de l'application.

\begin{itemize}[label=\ding{220}]
\item \textbf{La couche présentation (Front-End) :} cette couche gère l'ensemble des
interactions avec les utilisateurs. Elle se décline en deux interfaces complémentaires :
\begin{itemize}[label=\ding{43}]
    \item \textbf{React.js (application web)} : interface principale destinée aux
    administrateurs. Elle est structurée autour d'un routeur, de composants réutilisables
    et d'un service assurant la communication avec l'API via des requêtes HTTP ;
    \item \textbf{React Native / Expo (application mobile)} : interface mobile permettant
    la consultation du tableau de bord, la gestion des alertes et l'exécution des actions
    à distance en situation de mobilité.
\end{itemize}

\item \textbf{La couche métier (Back-End) :} implémentée en Node.js/Express, cette
couche assure le traitement des données, la logique métier et l'exposition des services
via une API REST. Elle est structurée en :
\begin{itemize}[label=\ding{43}]
    \item \textbf{Routes} : points d'entrée des requêtes HTTP, organisés par domaine
    fonctionnel (serveurs, métriques, alertes, actions à distance, sauvegardes) ;
    \item \textbf{Services} : implémentation de la logique fonctionnelle (calcul de
    statut, génération d'alertes, envoi d'emails, planification des sauvegardes) ;
    \item \textbf{Modèles} : schémas Mongoose assurant l'interaction avec la base de
    données.
\end{itemize}

\item \textbf{La couche sécurité :} assurée par un système d'authentification basé sur
les jetons JWT pour la gestion des sessions, complété par un middleware de vérification
du rôle administrateur protégeant l'ensemble des routes sensibles de l'API.

\item \textbf{La couche données :} repose sur MongoDB Atlas pour le stockage et la
persistance de l'ensemble des données de l'application (serveurs, métriques, alertes,
sauvegardes, journaux d'audit et utilisateurs).

\item \textbf{La couche collecte :} constituée des agents Python déployés sur chaque
serveur supervisé, transmettant périodiquement leurs métriques à la couche métier par
requêtes HTTP.

\item \textbf{La couche administration :} matérialisée par les connexions SSH établies à
la demande par le backend vers les serveurs supervisés, pour l'exécution des actions à
distance.
\end{itemize}

\begin{figure}[H]
    \centering
    \includegraphics[width=0.9\textwidth]{Image/ArchitectureLogique.png}
    \caption{Architecture logique de la plateforme}
    \label{fig:archi-logique}
\end{figure}

\subsubsection{Architecture physique de la plateforme}
L'architecture physique décrit la manière dont les différents composants logiciels sont
déployés et communiquent entre eux. La plateforme est déployée sur un serveur VPS OVH
sous Debian Linux, hébergeant un cluster Kubernetes allégé (k3s) qui centralise les
différents composants de l'application.

\begin{itemize}[label=\ding{220}]
\item \textbf{Frontend React.js :} l'application web est conteneurisée et servie par
Nginx au sein d'un pod Kubernetes. Elle est accessible depuis l'extérieur du cluster via
un service de type \texttt{NodePort} sur le port 30080.

\item \textbf{Backend Node.js/Express :} le serveur applicatif est exécuté dans un pod
Kubernetes dédié. Il expose l'API REST ainsi que le canal WebSocket via un service
\texttt{NodePort} sur le port 30300, et assure le traitement de la logique métier.

\item \textbf{MongoDB Atlas :} la base de données est hébergée dans le cloud sous forme
de service managé. Le backend y accède via une connexion sécurisée, les identifiants
étant injectés dans le pod au moyen d'un objet \texttt{Secret} Kubernetes.

\item \textbf{Agents Python :} déployés directement sur chaque serveur supervisé (VPS
OVH et machine virtuelle), ils transmettent leurs métriques au backend par requêtes HTTP
périodiques.

\item \textbf{Application mobile :} exécutée sur le terminal de l'utilisateur, elle
communique avec le backend via la même API REST que le client web.
\end{itemize}

\begin{figure}[H]
    \centering
    \includegraphics[width=0.9\textwidth]{Image/ArchitecturePhysique.png}
    \caption{Architecture physique de la plateforme}
    \label{fig:archi-physique}
\end{figure}

\section{Pilotage du projet avec Scrum}
Notre projet est piloté par Scrum, étant la méthode la plus sollicitée parmi les
méthodes agiles existantes.

\subsection{Backlog du produit}
Le tableau~\ref{tab:backlog-produit} présente le backlog du produit. Les fonctionnalités
sont organisées par module et décrites sous forme de \textit{User Stories}.

\renewcommand{\arraystretch}{1.5}
\begin{longtable}{|c|>{\raggedright\arraybackslash}p{2.6cm}|c|>{\raggedright\arraybackslash}p{6.2cm}|c|}
\hline
\textbf{ID} & \textbf{Module} & \textbf{ID US} & \centering\arraybackslash\textbf{User Story} & \textbf{Complexité} \\
\hline
\endfirsthead

\hline
\textbf{ID} & \textbf{Module} & \textbf{ID US} & \centering\arraybackslash\textbf{User Story} & \textbf{Complexité} \\
\hline
\endhead

\hline
\multicolumn{5}{r}{\textit{Suite du tableau à la page suivante}} \\
\endfoot

\hline
\caption{Backlog du produit}
\label{tab:backlog-produit}
\endlastfoot

M1 & Authentification
& M1.1 & En tant qu'administrateur, je souhaite me connecter à la plateforme avec mon email et mon mot de passe. & Moyenne \\
& & M1.2 & En tant qu'administrateur, je souhaite rester connecté grâce à un jeton persistant. & Moyenne \\
& & M1.3 & En tant qu'administrateur, je souhaite me déconnecter du système. & Faible \\
\hline

M2 & Gestion des serveurs
& M2.1 & En tant qu'administrateur, je souhaite consulter la liste des serveurs supervisés. & Moyenne \\
& & M2.2 & En tant qu'administrateur, je souhaite visualiser le statut de chaque serveur (OK, WARNING, CRITICAL, OFFLINE). & Élevée \\
& & M2.3 & En tant qu'administrateur, je souhaite consulter le détail d'un serveur. & Moyenne \\
& & M2.4 & En tant qu'administrateur, je souhaite qu'un nouveau serveur soit enregistré automatiquement dès sa première remontée de métriques. & Élevée \\
\hline

M3 & Supervision des métriques
& M3.1 & En tant qu'agent, je souhaite collecter les métriques système toutes les 5 secondes. & Élevée \\
& & M3.2 & En tant qu'administrateur, je souhaite visualiser les métriques en temps réel. & Élevée \\
& & M3.3 & En tant qu'administrateur, je souhaite consulter l'historique des métriques sous forme de graphique. & Élevée \\
& & M3.4 & En tant qu'administrateur, je souhaite consulter un résumé global de la santé du parc. & Moyenne \\
\hline

M4 & Gestion des alertes
& M4.1 & En tant que système, je souhaite générer une alerte lors du dépassement d'un seuil configuré. & Élevée \\
& & M4.2 & En tant qu'administrateur, je souhaite recevoir un email automatique lors d'une alerte. & Élevée \\
& & M4.3 & En tant qu'administrateur, je souhaite consulter la liste des alertes actives. & Moyenne \\
& & M4.4 & En tant qu'administrateur, je souhaite acquitter une alerte prise en charge. & Moyenne \\
\hline

M5 & Actions à distance
& M5.1 & En tant qu'agent, je souhaite détecter automatiquement les services actifs du serveur. & Élevée \\
& & M5.2 & En tant qu'administrateur, je souhaite consulter les services réellement actifs d'un serveur. & Moyenne \\
& & M5.3 & En tant qu'administrateur, je souhaite redémarrer un service à distance via SSH. & Élevée \\
& & M5.4 & En tant qu'administrateur, je souhaite arrêter un service à distance via SSH. & Élevée \\
& & M5.5 & En tant qu'administrateur, je souhaite redémarrer complètement un serveur à distance. & Élevée \\
& & M5.6 & En tant qu'administrateur, je souhaite consulter le journal d'audit des actions effectuées. & Moyenne \\
\hline

M6 & Gestion des sauvegardes
& M6.1 & En tant que système, je souhaite exécuter une sauvegarde automatique quotidienne. & Élevée \\
& & M6.2 & En tant qu'administrateur, je souhaite recevoir un email de confirmation après chaque sauvegarde. & Moyenne \\
& & M6.3 & En tant qu'administrateur, je souhaite consulter l'historique des sauvegardes. & Moyenne \\
\hline

M7 & Application mobile
& M7.1 & En tant qu'administrateur, je souhaite me connecter depuis l'application mobile. & Moyenne \\
& & M7.2 & En tant qu'administrateur, je souhaite consulter le tableau de bord depuis mon téléphone. & Élevée \\
& & M7.3 & En tant qu'administrateur, je souhaite exécuter des actions à distance depuis mon téléphone. & Élevée \\
\hline

M8 & Chaîne DevSecOps
& M8.1 & En tant que responsable DevOps, je souhaite conteneuriser chaque composant avec Docker. & Élevée \\
& & M8.2 & En tant que responsable DevOps, je souhaite automatiser le build via un pipeline Jenkins. & Élevée \\
& & M8.3 & En tant que responsable DevOps, je souhaite analyser la qualité du code avec SonarQube. & Moyenne \\
& & M8.4 & En tant que responsable DevOps, je souhaite scanner les vulnérabilités des images avec Trivy. & Moyenne \\
& & M8.5 & En tant que responsable DevOps, je souhaite déployer la plateforme sur un cluster Kubernetes. & Élevée \\
& & M8.6 & En tant que responsable DevOps, je souhaite automatiser le déploiement via GitOps (Argo CD). & Élevée \\
& & M8.7 & En tant que responsable DevOps, je souhaite sécuriser les secrets avec HashiCorp Vault. & Élevée \\
& & M8.8 & En tant que responsable DevOps, je souhaite superviser l'infrastructure avec Prometheus et Grafana. & Moyenne \\
\hline

\end{longtable}

\subsection{Planification des sprints}
Après l'élaboration du backlog produit et l'estimation de la charge de travail, le
projet a été réparti en quatre sprints. Chaque sprint regroupe un ensemble cohérent de
fonctionnalités permettant de livrer progressivement les différents modules de la
plateforme.

\begin{longtable}{|p{2.4cm}|>{\raggedright\arraybackslash}p{11.8cm}|}
\hline
\textbf{Sprint} & \textbf{Modules développés} \\
\hline
\endfirsthead

\hline
\textbf{Sprint} & \textbf{Modules développés} \\
\hline
\endhead

\hline
\caption{Planification des sprints}
\label{tab:planification-sprints}
\endlastfoot

\textbf{Sprint 1} &
\begin{itemize}[leftmargin=*, nosep]
  \item Authentification (M1)
  \item Gestion des serveurs (M2)
  \item Supervision des métriques (M3)
\end{itemize} \\
\hline

\textbf{Sprint 2} &
\begin{itemize}[leftmargin=*, nosep]
  \item Gestion des alertes (M4)
  \item Actions à distance (M5)
  \item Gestion des sauvegardes (M6)
  \item Application mobile (M7)
\end{itemize} \\
\hline

\textbf{Sprint 3} &
\begin{itemize}[leftmargin=*, nosep]
  \item Conteneurisation Docker (M8.1)
  \item Pipeline CI/CD Jenkins (M8.2)
  \item Analyse qualité SonarQube (M8.3)
  \item Scan de vulnérabilités Trivy (M8.4)
\end{itemize} \\
\hline

\textbf{Sprint 4} &
\begin{itemize}[leftmargin=*, nosep]
  \item Déploiement Kubernetes k3s (M8.5)
  \item GitOps avec Argo CD (M8.6)
  \item Gestion des secrets Vault et RBAC (M8.7)
  \item Supervision Prometheus et Grafana (M8.8)
\end{itemize} \\
\hline

\end{longtable}

\section{Environnement de travail}

\subsection{Environnement matériel de développement}
Le développement de l'application a été réalisé dans un environnement matériel adapté
pour répondre aux exigences du projet. Le poste de travail utilisé dispose des
caractéristiques suivantes :

\begin{itemize}[label=\ding{43}]
    \item \textbf{Modèle :} Ordinateur portable Lenovo
    \item \textbf{Système d'exploitation :} Windows 11 Pro (64 bits)
    \item \textbf{Processeur :} Intel(R) Core(TM) i7
    \item \textbf{Mémoire RAM :} 16 Go
    \item \textbf{Stockage :} 512 Go SSD
\end{itemize}

Le déploiement en production s'appuie quant à lui sur l'infrastructure suivante :

\begin{itemize}[label=\ding{43}]
    \item \textbf{Serveur VPS OVH :} Debian Linux, hébergeant le cluster Kubernetes k3s ;
    \item \textbf{Machine virtuelle VirtualBox :} second serveur supervisé permettant de
          valider le fonctionnement multi-serveurs ;
    \item \textbf{Smartphone Android :} utilisé pour les tests de l'application mobile.
\end{itemize}

\subsection{Outils de modélisation, de développement et de gestion de données}
Pour mener à bien la conception et le développement de la plateforme, plusieurs outils
ont été utilisés, allant de la modélisation UML à la documentation des API, en passant
par le développement intégré et la gestion des bases de données.

\vspace{0.5em}
\renewcommand{\arraystretch}{1.3}
\begin{center}
\begin{tabular}{|l|p{8.5cm}|c|}
\hline
\textbf{Outil} & \textbf{Description} & \textbf{Logo} \\
\hline
Draw.io & Logiciel en ligne utilisé pour la création des diagrammes UML (cas
d'utilisation, classes, séquences) et des schémas d'architecture de la plateforme.
& \includegraphics[height=25pt]{Image/drawio.png} \\
\hline
Visual Studio Code & Éditeur de code principal utilisé pour le développement du frontend
React.js, du backend Node.js/Express, de l'agent Python et de l'application mobile.
& \includegraphics[height=25pt]{Image/vscode.png} \\
\hline
Postman & Outil utilisé pour tester, documenter et valider les endpoints de l'API REST
de la plateforme (serveurs, métriques, alertes, actions à distance).
& \includegraphics[height=25pt]{Image/postman.png} \\
\hline
MongoDB Compass & Interface graphique officielle de MongoDB permettant de visualiser,
explorer et manipuler les données stockées dans la base de données.
& \includegraphics[height=25pt]{Image/compass.png} \\
\hline
GitHub & Plateforme de gestion de versions utilisée pour le suivi du code source et
constituant la source de vérité du déploiement GitOps.
& \includegraphics[height=25pt]{Image/github.png} \\
\hline
\end{tabular}
\captionof{table}{Outils utilisés pour la modélisation, le développement et la gestion de données}
\label{tab:outils}
\end{center}

\vspace{1em}

\subsection{Langages de programmation}
Le choix des langages de programmation est basé sur les exigences techniques de la
plateforme, notamment la performance du backend, la réactivité du frontend web et la
portabilité de l'application mobile.

\vspace{0.5em}
\renewcommand{\arraystretch}{1.3}
\begin{center}
\begin{tabular}{|l|p{8.5cm}|c|}
\hline
\textbf{Langage} & \textbf{Description} & \textbf{Logo} \\
\hline
JavaScript & Langage principal utilisé pour le développement du frontend (React.js) et
du backend (Node.js/Express). Il permet une cohérence technologique entre les deux
couches grâce à l'écosystème npm.
& \includegraphics[height=25pt]{Image/javascript.png} \\
\hline
TypeScript & Sur-ensemble de JavaScript intégrant des types statiques, utilisé pour le
développement de l'application mobile afin de sécuriser le code et faciliter sa
maintenance.
& \includegraphics[height=25pt]{Image/typescript.png} \\
\hline
Python & Utilisé pour développer l'agent de supervision installé sur chaque serveur. Sa
richesse en bibliothèques système, notamment \texttt{psutil}, en fait un choix idéal pour
la collecte de métriques.
& \includegraphics[height=25pt]{Image/python.png} \\
\hline
\end{tabular}
\captionof{table}{Langages de programmation utilisés dans la plateforme}
\label{tab:langages}
\end{center}

\vspace{1em}

\subsection{Frameworks et bibliothèques}
Les frameworks et bibliothèques retenus ont été sélectionnés pour leur maturité, leur
performance et leur adéquation aux besoins fonctionnels de la plateforme.

\vspace{0.5em}
\renewcommand{\arraystretch}{1.3}
\begin{center}
\begin{tabular}{|l|p{8.5cm}|c|}
\hline
\textbf{Framework} & \textbf{Description} & \textbf{Logo} \\
\hline
React.js & Bibliothèque JavaScript développée par Meta, utilisée pour construire
l'interface web. Son architecture basée sur les composants réutilisables garantit une
interface moderne et réactive.
& \includegraphics[height=25pt]{Image/react.png} \\
\hline
Node.js / Express & Framework backend JavaScript permettant de construire une API REST
performante et asynchrone. Il gère l'ensemble de la logique métier de la plateforme.
& \includegraphics[height=25pt]{Image/nodejs.png} \\
\hline
React Native / Expo & Framework permettant de créer l'application mobile à partir d'une
seule base de code, déployable sur Android et iOS.
& \includegraphics[height=25pt]{Image/reactnative.png} \\
\hline
Socket.io & Bibliothèque assurant la communication bidirectionnelle en temps réel entre
le serveur et les clients via le protocole WebSocket.
& \includegraphics[height=25pt]{Image/socketio.png} \\
\hline
\end{tabular}
\captionof{table}{Frameworks et bibliothèques utilisés dans la plateforme}
\label{tab:frameworks}
\end{center}

\vspace{1em}

\subsection{Outils DevOps et de déploiement}
La mise en production de la plateforme repose sur une chaîne DevSecOps automatisée
garantissant la fiabilité, la sécurité et la répétabilité des déploiements.

\vspace{0.5em}
\renewcommand{\arraystretch}{1.3}
\begin{center}
\begin{tabular}{|l|p{8.5cm}|c|}
\hline
\textbf{Outil} & \textbf{Description} & \textbf{Logo} \\
\hline
Docker & Outil de conteneurisation utilisé pour empaqueter l'application et ses
dépendances dans des conteneurs isolés, garantissant un environnement identique entre le
développement et la production.
& \includegraphics[height=25pt]{Image/docker.png} \\
\hline
Jenkins & Serveur d'intégration et de déploiement continu (CI/CD) utilisé pour
automatiser les pipelines de build, de test et de déploiement de la plateforme.
& \includegraphics[height=25pt]{Image/jenkins.png} \\
\hline
SonarQube & Plateforme d'analyse statique du code source, permettant de détecter les
\textit{code smells}, la duplication et les vulnérabilités avant tout déploiement.
& \includegraphics[height=25pt]{Image/sonarqube.png} \\
\hline
Trivy & Scanner de vulnérabilités open source spécialisé dans l'analyse des images de
conteneurs, comparant leurs composants aux bases de données de CVE connues.
& \includegraphics[height=25pt]{Image/trivy.png} \\
\hline
Kubernetes (k3s) & Orchestrateur de conteneurs allégé assurant le déploiement, la
mise à l'échelle et le redémarrage automatique des composants de la plateforme.
& \includegraphics[height=25pt]{Image/k3s.png} \\
\hline
Argo CD & Outil de déploiement continu selon l'approche GitOps, synchronisant
automatiquement l'état du cluster avec les manifestes versionnés dans Git.
& \includegraphics[height=25pt]{Image/argocd.png} \\
\hline
HashiCorp Vault & Solution de gestion centralisée des secrets, assurant le stockage
sécurisé et la distribution contrôlée des informations sensibles.
& \includegraphics[height=25pt]{Image/vault.png} \\
\hline
Prometheus / Grafana & Couple d'outils assurant la collecte des métriques
d'infrastructure et leur visualisation sous forme de tableaux de bord.
& \includegraphics[height=25pt]{Image/prometheus.png} \\
\hline
\end{tabular}
\captionof{table}{Outils DevOps et de déploiement de la plateforme}
\label{tab:devops}
\end{center}

\section*{Conclusion}
Nous nous sommes concentrés dans ce chapitre sur l'étude globale du projet afin de bien
le planifier. Nous avons identifié les acteurs, spécifié les besoins fonctionnels et non
fonctionnels, élaboré le backlog produit et planifié les sprints. Nous avons également
justifié le choix de l'architecture N-Tiers et présenté l'environnement de travail
retenu. Le chapitre suivant sera consacré à la réalisation du premier sprint.
